%    Latex Template for AMMCS
%%%%%%%%%%%%%%%%%%%%%%%%%%%%%%%%%%%%%%%%
%    DO NOT CHANGE THE SETTINGS BELOW
%%%%%%%%%%%%%%%%%%%%%%%%%%%%%%%%%%%%%%%%
\documentclass[11pt]{article}
\usepackage{epsfig}
\usepackage{graphicx}
\usepackage[latin1]{inputenc}
\usepackage[T1]{fontenc}
\usepackage{mathptmx}
\usepackage{url}
\setlength{\textheight}{24cm}		
\setlength{\textwidth}{17.5cm}		
\setlength{\oddsidemargin}{0.0cm}
\setlength{\parskip}{0.2cm plus 0.3cm minus 0.1cm}
\setlength{\topmargin}{-2.25cm}
\pagestyle{empty}
\begin{document}
\thispagestyle{empty}
\begin{tabbing}
\hspace*{4.5in} \= {} \\[-5pt]
\rule{6.8 in}{0.01 in} \\[-2pt]
{} \> {\footnotesize\tt  The VII International AMMCS} \\[-5pt]
{} \> {\footnotesize\tt  Interdisciplinary Conference} \\[-5pt]
{} \> {\footnotesize\tt  Waterloo, Ontario, Canada} \\[-5pt]
\end{tabbing}
%%%%%%%%%%%%%%%%%%%%%%%%
%%%%%%%%%%%%%%%%%%%%%%%%%%%%%%%%%%%%%%%%
%    Except the date,
%    DO NOT CHANGE THE SETTINGS ABOVE
%%%%%%%%%%%%%%%%%%%%%%%%%%%%%%%%%%%%%%%%
%
\vspace{-0.8cm}
\begin{flushleft}
\Large \textbf{\noindent
An Interpretable Transformer-based ABSA Framework for Validity-Weighted Attenuation of Affective Noise in Student Evaluations of Teaching}
\\
\vspace{0.3cm}
\normalsize
\normalsize{
\underline{B.~Obeidi}$^1$, S.~Fatahi$^1$, A. A. Bidgoli$^1$
} \\
\vspace{3mm}
\textit{\footnotesize
 $^1$ Wilfrid Laurier University, Ontario, Canada, \{obei4210@mylaurier.ca, sfatahi@wlu.ca, aasilianbidgoli@wlu.ca\}\\
}
\end{flushleft}

Student Evaluations of Teaching (SET) inform consequential institutional decisions on
faculty performance, tenure, and promotion, yet their validity is systematically compromised by well-documented biases~\cite{MR97}. In particular, non-pedagogical content distorts numerical ratings in ways unrelated to instructional quality~\cite{W11}. Informal platforms
such as RateMyProfessors.com (RMP) exacerbate this through full anonymity and absence of
moderation. Despite advances in computational SET analysis, existing work remains
descriptive: sentiments are extracted but ratings are never intervened upon. This work
addresses that gap with a five-stage interpretable Aspect-Based Sentiment Analysis (ABSA)
framework that directly adjusts SET ratings to attenuate affective noise while preserving
pedagogical signal and students' evaluative intent.

The framework is applied to $N = 17{,}121$ English-language RMP reviews. Clauses are
extracted via spaCy and assigned to one of three pedagogical topics (\textit{Instructional
Effectiveness}, \textit{Fairness \& Grading}, \textit{Workload \& Difficulty}) or to a
\textit{Miscellaneous} category via zero-shot Sentence-BERT cosine similarity thresholding.
A pedagogical density metric ($D_\mathrm{misc}$) captures the proportion of non-pedagogical content per review. Clause-level emotion vectors
(27 non-neutral dimensions) are extracted via RoBERTa-GoEmotions and aggregated at the topic level.
A CatBoost regressor trained on these features predicts ratings; adjusted scores are produced by recomputing predictions under proportionally attenuated \textit{Miscellaneous} emotional inputs via attenuation factor
$\zeta = 1 - (D_{\mathrm{misc}})^{s}$ and computing
\begin{equation}
  \hat{y}^{\mathrm{adjusted}} = y^{\mathrm{raw}} + \bigl(\hat{y}^{\mathrm{attenuated}} - \hat{y}^{\mathrm{baseline}}\bigr). \label{eq1}
\end{equation}
This prediction-based correction mechanism aims to keep adjustments attributable to specific review text and topic-level emotions, satisfying interpretability requirements for consequential decision-making~\cite{DK17}.

The regressor achieves MAE $= 0.637$, Spearman $\rho = 0.736$, $R^2 = 0.614$ on a
professor-level stratified hold-out set. Aspect-Term Categorization achieves an expert-validated accuracy of 0.77. SHAP-based internal validation
confirms attenuation shifts model sensitivity away from \textit{Miscellaneous} features ($-10.7\%$) toward
pedagogical categories (\textit{Instructional Effectiveness}: $+2.1\%$; \textit{Fairness}: $+2.7\%$; \textit{Workload}:
$+2.6\%$), while maintaining or strengthening correlations with pedagogical features.
External validation with a senior faculty expert yielded paired comparison accuracy of
$88.9\%$ under coarse conditions ($p < 0.001$) and $75.9\%$ for fine-grained ($p < 0.005$).

To our knowledge, this is the first transparent computational mechanism to operationalize
construct-irrelevant variance (CIV) mitigation directly into SET rating adjustment. The framework bridges Messick's unified
validity theory~\cite{M95} with classical reliability and robust estimation principles,
yielding auditable corrections grounded in psychometric theory. The modular design also
suggests applicability to other open-ended feedback domains where reviews are accompanied by numerical ratings.

\vspace{-0.2cm}
%%%%%%%%%%%%%%%%%%%%%%%%%%%%%%%%%%%%%%%%
%% BIBLIOGRAPHY
%%%%%%%%%%%%%%%%%%%%%%%%%%%%%%%%%%%%%%%%
\begin{thebibliography}{00}
\addcontentsline{toc}{chapter}{References}
\setlength{\itemsep}{-1mm}
\footnotesize
\bibitem{MR97} H.~W.~Marsh and L.~A.~Roche, \textit{Making students' evaluations of teaching effectiveness effective}, Amer. Psychologist \textbf{52}(11), pp.~1187--1197 (1997).
\bibitem{W11} W.~Wongsurawat, \textit{What's a comment worth? How to better understand student evaluations of teaching}, Quality Assurance in Education \textbf{19}(1), pp.~67--83 (2011).
\bibitem{DK17} F.~Doshi-Velez and B.~Kim, \textit{Towards a rigorous science of interpretable machine learning}, arXiv:1702.08608 (2017).
\bibitem{M95} S.~Messick, \textit{Validity of psychological assessment}, Amer. Psychologist \textbf{50}(9), pp.~741--749 (1995).
\end{thebibliography}
\normalsize
\end{document}